\documentclass[12pt,a4paper]{report}
\usepackage[utf8]{inputenc}
\usepackage{amsmath}
\usepackage{amsfonts}
\usepackage{amssymb}
\usepackage{amsthm}
\usepackage{hyperref}

\usepackage{multicol}
\usepackage{fancyhdr}
\usepackage[inline]{enumitem}
\usepackage{tikz}
\usepackage{tikz-cd}
\usetikzlibrary{calc}
\usetikzlibrary{shapes.geometric}
\usetikzlibrary{positioning}
\usepackage[margin=0.5in]{geometry}
\usepackage{xcolor}

\hypersetup{
    colorlinks=true,
    linkcolor=blue,
    filecolor=magenta,      
    urlcolor=cyan,
    pdftitle={Tensors},
    pdfpagemode=FullScreen,
    }

%\urlstyle{same}

\newcommand{\CLASSNAME}{Math 5050 -- Special Topics: Tensor Analysis}
\newcommand{\STUDENTNAME}{Paul Carmody}
\newcommand{\ASSIGNMENT}{Chapter 9: Determinants and the Levi-Civita Symbol }
\newcommand{\DUEDATE}{February 2026}
\newcommand{\PROFESSOR}{Professor Berchenko-Kogan}
\newcommand{\SEMESTER}{Spring 2026}
\newcommand{\SCHEDULE}{TBD}
\newcommand{\ROOM}{Remote}

\newcommand{\MMN}{M_{m\times n}}
\newcommand{\FF}{\mathcal{F}}

\pagestyle{fancy}
\fancyhf{}
\chead{ \fancyplain{}{\CLASSNAME} }
%\chead{ \fancyplain{}{\STUDENTNAME} }
\rhead{\thepage}

\fancyfoot{} % clear all footer fields
\fancyfoot[LE,RO]{\thepage}
\fancyfoot[LO,CE]{-------\\\ASSIGNMENT}
%\fancyfoot[CO,RE]{To: Dean A. Smith}
\input{../MyMacros}

\newcommand{\RED}[1]{\textcolor{red}{#1}}
\newcommand{\BLUE}[1]{\textcolor{blue}{#1}}

\newcommand{\SUPP}{\operatorname{supp}}
\newcommand{\CLOSURE}{\operatorname{closure of}}
\newcommand{\BD}{\operatorname{bd}}
\newcommand{\HHN}{\mathcal{H}^n}
\newcommand{\COFF}[3]{\Gamma^{#1}_{{#2}{#3}}}
\newcommand{\VK}[1]{\textbf{#1}}
\newcommand{\VKR}{\VK{R}}
\newcommand{\VKZ}{\VK{Z}}

\begin{document}

\begin{center}
	\Large{\CLASSNAME -- \SEMESTER} \\
	\large{ w/\PROFESSOR}
\end{center}
\begin{center}
	\STUDENTNAME \\
	\ASSIGNMENT -- \DUEDATE\\
\end{center} 

\textbf{\large{Note:}}
\begin{description}
	\item \textbf{Cross Product}
	
	We define the Cross Product as 
	\begin{align*}
		W = U \times V \implies	W^i= e^{ijk}U_jV_k
	\end{align*}That is 
	\begin{align*}
		\forall i,\, W^i = \sum_j\sum_k e^{ijk}U_jV_k
	\end{align*}That is, $e^{ijk}$ sets the sign for the term, adding when $ijk$ are in order and subtracting when they are in reverse order.  Thus, the vector equation becomes
	\begin{align*}
		\VK{W} &= e^{ijk}U_jV_k\,\VK{Z}_i
	\end{align*}
\end{description}
\HLINE
\noindent\textbf{Exercises: 9.2}
\begin{description}

\item \textbf{Exercise 161}  Show that equation (9.3) follows from equation (9.1) and (9.2).
\begin{align*}
	A_{ijk} &= -A_{jik} & (9.1) \\
	A_{ijk} &= -A_{ikj} & (9.2) \\
	A_{ijk} &= -A_{kji} & (9.3)
\end{align*}

\BLUE{From 9.1 $\implies$ 9.3: flip twice 1) $k \leftrightarrow j = kij$ and 2) $j \leftrightarrow i = kji$.  Each flip will change the sign. \\
From 9.2 $\implies$ 9.3:  flip twice 1) $i \leftrightarrow j = kij$ and 2) $i \leftrightarrow j = kji$.  Each flip will change the sign.
}

\item \textbf{Exercise 162} Evaluate the contraction $e^{ijk}e_{ijk}$.

\BLUE{\begin{align*}
	e^{ijk}e_{ijk} &= \sum_i\sum_j\sum_k e^{ijk}e_{ijk} 
\end{align*}Observe when $i=1$ we have 
\begin{align*}
	e^{1jk}e_{1jk} =& e^{111}e_{111}+e^{112}e_{112} + e^{113}e_{113} + \\
	&\text{ }e^{121}e_{121}+e^{122}e_{122} + e^{i23}e_{123} +\\
	& e^{131}e_{131}+e^{132}e_{132} + e^{i33}e_{133} \\
	=& 0+0 + 0 + \\
	& 0+0 + e^{123}e_{123} +\\
	& 0+e^{132}e_{132} + 0
\end{align*}The zeros occur whenever indices appear twice.  Then,
\begin{align*}
	e^{ijk}e_{ijk} = e^{123}e_{123} + e^{132}e_{132} &= 1 + (-1)(-1) = 2\end{align*}This is true for $i=2$ and $i=3$.  Thus, $e^{ijk}e_{ijk} = 6$
}
\item \textbf{Exercise 163}  Evaluate $e^{ijk}e_{kij}$.

\BLUE{This is the same argument as 162.  Observe that $i=1$ the order of the indices doesn't matter.  If they are duplicated \textit{either} above or below (in this case both) the term becomes 6.
}

\item \textbf{Exercise 164}  What are the entries of the sysmem $f_{ijk}=e_{ijk}+e_{jik}$.

\BLUE{Ignoring any repeated indices (which would be zero), for any $ijk$ we get a value but the $e_{jik}$ will be the negative of the same value.  Thus, $f_{ijk}=0$ for all values.
}

\item \textbf{Exercise 165} What are the entries of the system $f_{ijk}=e_{ijk}+e_{kij}$?

\BLUE{Ignoring any repeated indices (which would be zero), we can see that $kij$ is in the same cyclic order as $ijk$ thus. 
\begin{align*}
	f_{ijk} &= \BINDEF{2 & ijk \text{ are in cyclic order}}{0 & \text{otherwise}}
\end{align*}
}

\HLINE
\noindent\textbf{Section 9.3}

\item \textbf{Exercise 166.} Explain each step in equation (9.11)
\begin{align*}
	e^{ijk}a_i^2a_j^1a_k^3 \overset{A}= e^{ijk}a_j^1a_i^2a_k^3\overset{B}=e^{jik}a_i^1a_j^2a_k^3\overset{C}=-e^{ijk}a_i^1a_j^2a_k^3 \overset{D}= A.
\end{align*}
\BLUE{\begin{description}
	\item \textbf{A: }swap the positions of $a_i^2 \leftrightarrow a_j^1$.  This is merely multiplicative commutativity.
	\item \textbf{B: }swap the both lower of $a$ and upper indices of $e$, $i \leftrightarrow j$. The upper indices changes the sign, the lower indices changes it back.
	\item \textbf{C: }now both upper and lower indices of $a$ are in increasing order but the indices of $e$ are negative.  Switch this to the direct cyclic form and reverse the sign.
	\item \textbf{D: }this is from the definition of determinant.
	\begin{align*}
		A &= \frac{1}{3!}e^{ijk}e^{rst}a_{ir}a_{js}a_{kt} & (9.15) \\
		A &= \frac{1}{3!}e_{ijk}e_{rst}a^{ir}a^{js}a^{kt} & (9.15)
	\end{align*}
\end{description}
}
\item \textbf{Exercise 167.} Derive equations (9.15) and (9.16).

\section{Section 9.4}
\item \textbf{Exercise 168} Justify equation (9.20)

\begin{align*}
	\delta_{rst}^{ijk} &= \left |\begin{array}{ccc}
		\delta_r^i & \delta_r^j & \delta_r^k \\
		\delta_s^i & \delta_s^j & \delta_s^k \\
		\delta_t^i & \delta_t^j & \delta_t^k \\
	\end{array} \right | & (9.20)
\end{align*}


\BLUE{This is only non-zero when $ijk$ and $rst$ are each permutations of $1, 2,3$.  Which gives us $2 * 3!$ possible combinations.  Half of these will have the same sign or permutation and the other will have opposite sign.\\
Construct the $3 \time 3$ matrix and it will become clear that there can only be one entry in each row and column.  Thus, when $a_{11}=1$ then there are two options: 1) is that $a_{22}=1$ which implies that both permutations are the same and the detrminant is 1  or 2) is that $a_{23}=1$ which indicates that the permuatations are not the same and the determinant is -1.\\
Thus, when $a_{12}=1$ we start by multiplying by -1 and, once again we have two choices that put the permutations together or opposite.\\
Similarly, when $a_{13}=1$.
}

\item \textbf{Exercise 169} Explain why equation (9.20) shows that $\delta_{rst}^{ijk}$ is a tensor.

\BLUE{(9.20) is the sum of products of tensors.  Each Kronecker Delta is a tensor and the determinant is built from terms each of which is a tensor.
}
\item \textbf{Exercise 170} Explain why it follows that the determinant of a tensor with one covariant and one contravariant index is invariant.
\item \textbf{Exercise 171} Justify equations (9.21) and (9.25).

\begin{align*}
	\delta_{rs}^{ij} &= \delta_{rsk}^{ijk} & (9.21)\\
	3\delta_j^i &= \delta_{rj}^{ij} & (9.25)
\end{align*}

\BLUE{\begin{align*}
	\delta_{rsk}^{ijk} &= \left |\begin{array}{ccc}
		\delta_r^i & \delta_r^j & \delta_r^k \\
		\delta_s^i & \delta_s^j & \delta_s^k \\
		\delta_k^i & \delta_k^j & \delta_k^k \\
	\end{array} \right | \\
	&= \delta_k^i \DTWOXTWO{\delta_r^j}{\delta_r^k}{\delta_s^j}{\delta_s^k} - \delta_k^j\DTWOXTWO{\delta_r^i}{\delta_r^k}{\delta_s^i}{\delta_s^k} + \delta_k^k\DTWOXTWO{\delta_r^i}{\delta_r^j}{\delta_s^i}{\delta_s^j} \\
	&= \delta_k^i\delta_r^j\delta_s^k-\delta_k^i\delta_r^k\delta_s^j - \delta_k^j\delta_r^i\delta_s^k+\delta_k^j\delta_r^k\delta_s^i+\delta_k^k\delta_r^i\delta_s^j-\delta_k^k\delta_r^j\delta_s^i \\
	&= \delta_r^j\delta_s^i-\delta_r^i\delta_s^j-\delta_s^j\delta_r^i+\delta_r^j\delta_s^i+\delta_r^i\delta_s^j-\delta_r^j\delta_s^i \\
	&= \delta_r^j\delta_s^i-\delta_s^j\delta_r^i \\
	&= \delta_{rs}^{ij}
\end{align*}
}

\item \textbf{Exercise 172} Justify equations (9.23) and (9.25) in four dimensions.

\begin{align*}
	\delta_{rst}^{ijk} &= \delta_{rstl}^{ijkl} & (9.23)\\
	3\delta_j^i &= \delta_{rj}^{ij} & (9.25)
\end{align*}

\item \textbf{Exercise 173} Generalize equations (9.23) and (9.25) to $N$ dimensions.
\item \textbf{Exercise 174} Justify equation (9.26)
\item \textbf{Exercise 175} Evaluate $\delta_{ijk}^{ijk}$.

\BLUE{it equals zero when not a permutation.  Otherwise, it is always one.
}

\item \textbf{Exercise 176} Evaluate $\delta_{ikj}^{jir}$.
\item \textbf{Exercise 177} Evaluate $\delta_{ijk}^{rst}\delta_{lmn}^{rst}$.
\item \textbf{Exercise 178} Evaluate $\delta_{ijk}^{rst}\delta_{lmn}^{rst}\delta_{ijk}^{lmn}$.

\section{Section 9.5}
\item \textbf{Exercise 179} Show $C=AB$ property for systems $c_k^i, a^{ij},$ and $b_{jk}$ related by $c_k^i=a^{ij}b_{jk}$.

\section{Section 9.6}
\item \textbf{Exercise 180} Explain equation (9.39).
\item \textbf{Exercise 181} Show that each term in equation (9.41) equals $2A_l^u$.
\item \textbf{Exercias 182} Generalize the proof of equation (9.36) to arbitrary dimension.
\item \textbf{Exercise 183} Derive similar expressions for derminants of tensors $a_{ij}$ and $a^{ij}$.
\item \textbf{Exercise 184} Explain why each term vanishes on the right-hadn side of equation (9.46).
\item \textbf{Exercise 185} Explain why each term in parentheses in equation (9.48) equals $A$.
\item \textbf{Exercise 186} Extend the subject of cofactors to systems with two lower indices (denote the cofactor by $A^{ir}$) and systems with two upper indices (denote the cofactor by $A_{ir}$).

\section{Section 9.7}
\item \textbf{Exercise 187} Confirm equations (9.57), (9.58), and (9.59).
\item \textbf{Exercise 188} Use the Voss-Weyl formula to dericve the Laplacian in Spherical Coordinates.
\item \textbf{Exercise 189} Use the Voss-Weyl formula to derive the epxression (8.37) for the Laplacian in cylindrical coordinates

\section{Section 9.9}
\item \textbf{Exercise 190} Show that if $S_{jk}^i$ is a relative tensor of weight $M$ and $T_t^{rs}$ is a relative tensor of weight $N$, then $S_{jk}^iT_t^{rs}$ is a relative tensor of weight $M + N$.  In particular, if $M=-N$, then $S_{jk}^iT_t^{rs}$ is a tensor.

\item \textbf{Exercise 191} Show that the result of contraction for a relative tensor of weight $M$ is also a relative tensor of weight $M$.

\BLUE{\begin{align*}
	\CONJ{T}	_a^n &= J^MT_m^bJ_b^aJ_n^m
\end{align*}A contraction would be $a=n$ or
\begin{align*}
	\CONJ{T}	_a^a &= J^MT_m^bJ_b^aJ_a^m = J^MT_m^bJ_b^m
\end{align*}
}

\item \textbf{Exercise 192} Conclude that $\delta_{rst}^{ijk}$ is an absolute tensor of the basis in equation (9.17)

\item \textbf{Exercise 193} Show that the determinant of a relative covariant tensor $a_{ij}$ of weight $M$ is a relative invariant of weith $2+3M$ and tha tin $n$ dimensions the expression generalizes to $2+nM$.

\item \textbf{Exercise 194} Show that the volume element $\sqrt{Z}$ is a relative invariant of weight 1 with respect to orientation-preserving coordinate changes.

\section{Section 9.10}
\item \textbf{Exercise 195} Confirm that the Levi-Civita symbols are absolute tensors with respect to orientation preserving coordinate changes.

\begin{align*}
	\varepsilon^{ijk} &= \frac{e^{ijk}}{\sqrt{Z}} \\
	\varepsilon_{ijk} &= e_{ijk}\sqrt{Z}
\end{align*}

\BLUE{Recall that \begin{align*}
	e_{ijk}J_{i'}^iJ_{j'}^jJ_{k'}^k &= Je_{i''j'k} & (9.74) 
\end{align*} then by definition
\begin{align*}
	\varepsilon_{ijk}J_{i'}^iJ_{j'}^jJ_{k'}^k &= J_{i'}^iJ_{j'}^jJ_{k'}^ke_{ijk}\sqrt{Z} \\
	&= Je_{i'j'k'} \sqrt{Z'/J^2} \\
	&= Je_{i'j'k'} J^{-1}\sqrt{Z'} \\
	&= e_{i'j'k'} \sqrt{Z'}
\end{align*}
}

\item \textbf{Exercise 196} Show that 
\begin{align*}
	\delta_{rst}^{ijk} = \varepsilon^{ijk}\varepsilon_{rst},
\end{align*}which offers yet another proof that $\delta_{rst}^{ijk}$ is an absoltue tensor.
\item \textbf{Exercise 197} Show that te permutation $e^{ijk}$ and $e_{ijk}$ are \textbf{not} related by index juggling.
\section{Section 9.11}
\item \textbf{Exercise 198} Show that the metrinilic property applies to the contrvariant Levi-Civita symbol $\varepsilon^{rst}$.
\item \textbf{Exercise 199} Use the metrinilic property of the Levi-Civita symbols to show same for the delta systems:
\begin{align*}
	\nabla_m\delta_{rst}^{ijk}, \, \nabla_m\delta_{rs}^{ij}, \, \nabla_m \delta_r^i = 0
\end{align*}

\section{Section 9.12}
\item \textbf{Exerise 200}  Show that $\varepsilon^{ijk}U_iU_jV_k = 0$ by rename-switch-reschuffle
\item \textbf{Exercise 201} Show that $|\VK{W}|=|\VK{U}||\VK{V}|\sin \alpha$ where $\sin\alpha$ is the angle between $\VK{U}$ and $\VK{V}$.
\item \textbf{Exercise 202} As practie with working with Levi-Civita sybols, derivet eh expression for the corss product of three vectors.  You should obtain an expression with dyadic form that reads
\begin{align*}
	(\textbf{U}\times \textbf{V})\times \textbf{W} = (\textbf{U}\cdot\textbf{W})\textbf{V} - (\textbf{V}\cdot\textbf{W})\textbf{U}
\end{align*}Conclude that the cross product is not associative.
\item \textbf{Exercise 203} Show that the indices can be roated in the expression for $\textbf{W}$ to yield
\begin{align*}
	\textbf{W} &= \varepsilon^{ijk}U_iV_j\textbf{Z}_k.
\end{align*}
\item \textbf{Exercise 204} Show that \begin{align*}
	\textbf{U}\cdot(\textbf{V}times \textbf{W}) = (\textbf{U}\times\textbf{V})\cdot\textbf{W}.
\end{align*}
\item \textbf{Exercise 205} Show that in Cartesian coordinates, the cross product is given by 
\begin{align*}
	\textbf{W} &= \left | \begin{array}{ccc}
		\textbf{i} & \textbf{j} & \textbf{k}\\
		U^1 & U^2 & U^3\\
		V^1 & V^2 & V^3
\end{array}	 \right |
\end{align*}
\item \textbf{Exercis 206} Show that in cylindrical coordinates, the cross product is give by  
\begin{align*}
	\textbf{W} &= \left | \begin{array}{ccc}
		r\textbf{Z}^1 & r\textbf{Z}^2 & r\textbf{Z}^3\\
		U^1 & U^2 & U^3\\
		V^1 & V^2 & V^3
\end{array}	 \right |
\end{align*}
\item \textbf{Exercise 207} Derive the epxression for the corss product in spherical coordinates.
\section{Section 9.13}
\item \textbf{Exercise 208} Analyze the curl applied to a cross product.  Show that $\textbf{W}= \nabla\times(\textbf{U}\times\textbf{V})$ is given by 
\begin{align*}
	W^i = -\nabla_jU^jV^i+\nabla_jU^iV^j-U^j\nabla_jV^i+U^i\nabla^jV^j.
\end{align*}
\item \textbf{Exercise 209} Show that divergence of curl vanishes.
\item \textbf{Exercise 210} Show that curl of a gradient vanishes.
\section{Section 9.14}
\item \textbf{Exercise 211} Expplain why $\varepsilon_{ijk}U^iV^i = 0$ in equation (9.135).
\item \textbf{Exercise 212} Show that the length of the vector $V^i$ equals the length of the vector $U^i$
\end{description}
\end{document}
